\uuid{LdT3}
\exo7id{777}
\titre{exo7 777}
\auteur{ridde}
\organisation{exo7}
\datecreate{1999-11-01}
\isIndication{false}
\isCorrection{false}
\chapitre{Fonctions circulaires et hyperboliques inverses}
\sousChapitre{Fonctions hyperboliques et hyperboliques inverses}
\cpge{mpsi-mp2i.fonctions-valeurs-reelles-complexes.c, pcsi.techniques-calcul-fonctions.c, ptsi.techniques-calcul-fonctions.c, mpsi-mp2i.calcul-algebrique-trigonometrie.d, pcsi.calcul-algebrique-trigonometrie.d, ptsi.calcul-algebrique-trigonometrie.d, bcpst1.trigonometrie}
\module{Analyse}
\niveau{L1}
\difficulte{}

\contenu{
\question{
R\'esoudre l'\'equation $\tan (3 \arcsin x) = 1$. On exprimera les trois solutions
au moyen de radicaux.
}
}
